Below there is a list of all the tools, libraries and sources used to compose this document.
\begin{itemize}
    \item \underline{\textbf{\href{https://plantuml.com}{PlantUML}}}: a tool used to create UML diagrams (sequence diagrams, deployment diagram and component integration diagrams) by writing textual descriptions.
    \item \underline{\textbf{\href{https://miktex.org}{MiKTeX}}}: A distribution of LaTeX for Windows, used to compile and render the LaTeX code into a PDF document.
    \item \underline{\textbf{\href{https://code.visualstudio.com}{Visual Studio Code}}}: A source-code editor used for writing, editing, and managing the LaTeX document and related code.
    \item \underline{\textbf{\href{https://marketplace.visualstudio.com/items?itemName=James-Yu.latex-workshop}{LaTeX Workshop extension for VS Code}}}: An extension for Visual Studio Code that adds support for LaTeX documents.
    \item \underline{\textbf{\href{https://strawberryperl.com}{Strawberry PERL}}}: A Perl distribution required by LaTeX tools.
    \item \underline{\textbf{\href{https://github.com}{GitHub}}}: a platform used for version control and collaborative work, where the LaTeX files and other resources were stored and managed.
    \item \underline{\textbf{\href{https://draw.io}{Draw.io}}}: it is a free, user-friendly tool for creating diagrams (component diagram and overview diagram in \underline{\hyperref[section2]{section 2}})
    \item \underline{\textbf{\href{https://www.figma.com/}{Figma}}}: it is a design tool used for interface design, prototyping, and creating interactive user experiences (see mockups in \underline{\hyperref[section3]{section 3}})
\end{itemize}